%\documentclass[12pt,a4paper]{report}
\documentclass[12pt,a4paper,oneside,onecolumn,openright]{book}
% set the document language
\usepackage[italian]{babel}
% set the encoding used by your editor here (default is utf8)
\usepackage[utf8]{inputenc}
\usepackage[T1]{fontenc}

% math packages
\usepackage{amsmath}
\usepackage{amssymb}
% page margins settings
\usepackage[inner=3cm,outer=2.5cm,top=3cm,bottom=2.5cm]{geometry}
%\usepackage{indentfirst}

% other packages
\usepackage{array}
\usepackage{subfigure}
\usepackage{graphicx}
\usepackage{verbatim}
\usepackage{listings}
\usepackage{url}
\usepackage[hidelinks]{hyperref}
% custom colors
\usepackage{color}
\definecolor{light-gray}{gray}{0.96}
\definecolor{cyan}{RGB}{230,230,255}
\definecolor{dkgreen}{rgb}{0,0.6,0}
\definecolor{gray}{rgb}{0.5,0.5,0.5}
\definecolor{mauve}{rgb}{0.58,0,0.82}

% environment for bash code
\lstset{ %
  language=bash,                % the language of the code
  basicstyle=\footnotesize,           % the size of the fonts that are used for the code
  numbers=left,                   % where to put the line-numbers
  numberstyle=\footnotesize,          % the size of the fonts that are used for the line-numbers
  stepnumber=1,                   % the step between two line-numbers. If it's 1, each line 
                                  % will be numbered
  numbersep=5pt,                  % how far the line-numbers are from the code
  backgroundcolor=\color{white},      % choose the background color. You must add \usepackage{color}
  showspaces=false,               % show spaces adding particular underscores
  showstringspaces=false,         % underline spaces within strings
  showtabs=false,                 % show tabs within strings adding particular underscores
%  frame=single,                   % adds a frame around the code
  rulecolor=\color{black},        % if not set, the frame-color may be changed on line-breaks within not-black text (e.g. commens (green here))
  tabsize=2,                      % sets default tabsize to 2 spaces
  captionpos=b,                   % sets the caption-position to bottom
  breaklines=true,                % sets automatic line breaking
  breakatwhitespace=false,        % sets if automatic breaks should only happen at whitespace
  title=\lstname,                   % show the filename of files included with \lstinputlisting;
                                  % also try caption instead of title
  numberstyle=\tiny\color{gray},        % line number style
  keywordstyle=\textbf,          % keyword style
  commentstyle=\color{dkgreen},       % comment style
%  stringstyle=\color{mauve},         % string literal style
  escapeinside={\%*}{*)},            % if you want to add a comment within your code
  morekeywords={*,...,insert,-}               % if you want to add more keywords to the setù
}

% environment for python code
\lstset{
language=Python,
breaklines=true,
breakatwhitespace=true ,
backgroundcolor=\color{light-gray}
}
% appendices package
%\usepackage{appendix}
% set Appendix name used in the toc
%\renewcommand{\appendixtocname}{Appendice}

% interline
\linespread{1.5}
% set numbers for subsections and show them in the toc
\setcounter{tocdepth}{3} 
\setcounter{secnumdepth}{3}

% layout package, style and settings
\usepackage{fancyhdr}
\pagestyle{fancy}

\fancypagestyle{mainmatter}{%		
		\fancyhf{} 
		\fancyhead{}
		\fancyhead[LE,RO]{\thepage}
		\fancyhead[LO]{\footnotesize{\leftmark}}
		\fancyhead[RE]{\footnotesize{\rightmark}}
		\fancyfoot{}
		\addtolength{\headwidth}{\marginparsep}
		\addtolength{\headheight}{2.5pt}
		\renewcommand{\headrulewidth}{0.3pt}
		\renewcommand{\footrulewidth}{0.0pt}
		}
\fancypagestyle{frontmatter}{%
		\fancyhf{} 
		\fancyhead[LE]{\footnotesize{\MakeUppercase{\thepage}}}
		\fancyhead[RO]{\footnotesize{\MakeUppercase{\thepage}}}
		\fancyhead[RE,LO]{}
		\fancyfoot{}
		\addtolength{\headwidth}{\marginparsep}
		\addtolength{\headheight}{2.5pt}
		\renewcommand{\headrulewidth}{0.0pt}
		\renewcommand{\footrulewidth}{0.0pt}
		}
		
		
\usepackage{fancyhdr}
\pagestyle{fancy}
		\fancyhf{} 
		\fancyhead{}
		\fancyhead[LE,RO]{\thepage} 
		\fancyhead[LO]{\footnotesize{\leftmark}}
		\fancyhead[RE]{\footnotesize{\rightmark}}
		\fancyfoot{}
		\addtolength{\headwidth}{\marginparsep}
		\addtolength{\headheight}{2.5pt}
		\renewcommand{\headrulewidth}{0.3pt}
		\renewcommand{\footrulewidth}{0.0pt}

% empty pages have no numbers
\makeatletter
\def\cleardoublepage{\clearpage\if@twoside \ifodd\c@page\else
\hbox{}
  %Potresti voler togliere il commento dalla linea seguente
  %Questa pagina � stata lasciata intenzionalmente vuota.
\thispagestyle{empty}
\newpage
\if@twocolumn\hbox{}\newpage\fi\fi\fi}
\makeatother
%????
%\textwidth=450pt\oddsidemargin=0pt

%\makeatletter 
%  \DeclareRobustCommand*\textsubscript[1]{% 
%    \@textsubscript{\selectfont#1}} 
%  \newcommand{\@textsubscript}[1]{% 
%    {\m@th\ensuremath{_{\mbox{\fontsize\sf@size\z@#1}}}}} 
\makeatother 

\begin{document}
\begin{titlepage}
\begin{center}
{
    \large
    \textbf{Università  degli studi di Modena e Reggio Emilia} \\
   	\textbf{Dipartimento di Ingegneria "Enzo Ferrari"} \\
    \vspace{\stretch{0.5}}
    \hspace*{0cm} \hrulefill \hspace*{0cm} \\
    \vspace{\stretch{0.5}}
   	\emph{Corso di Laurea in Ingegneria Informatica - Sede di Mantova}
    
	  \vspace{\stretch{12}}
  
  
 		\huge{\bf Titolo: prima riga }}\\
		\vspace{3mm}
		{\huge{\bf Seconda riga}}\\
		\vspace{3mm}
		\vspace{3mm}
		{\huge{\bf Terza riga}}\\
		\vspace{3mm}
		\vspace{3mm}
		{\huge{\bf Quarta riga}}\\
		
		\vspace{\stretch{6}}
		\end{center}
		
\vspace{40mm}
\par
\noindent
\begin{minipage}[t]{0.47\textwidth}
{\large{\bf Relatore:\\
Prof. Luca Ferretti}}\\ 
\\
{\large{\bf Correlatore:\\
Ing. Federico Magnanini}}
\end{minipage}
\hfill
\begin{minipage}[t]{0.47\textwidth}\raggedleft
{\large{\bf Candidato:\\
Giulio Barabino}}
\end{minipage}
\vspace{20mm}
\begin{center}
%\rule[0.1cm]{15.8cm}{0.1mm}
\hspace*{0cm} \hrulefill \hspace*{0cm} \\
{\large{\bf 
Anno Accademico 2021/2022}}
\end{center}

\end{titlepage}

\pagestyle{frontmatter}
\frontmatter

% PAGINA VUOTA
%\clearpage\null\thispagestyle{empty}\clearpage
\setcounter{tocdepth}{2}
\tableofcontents

\setlength{\parindent}{12pt}
\setlength{\parskip}{1ex plus 0.5ex minus 0.2ex}
\mainmatter
\pagestyle{mainmatter}

\chapter{Introduzione}
\label{intro}

Il Single Sign-On è un protocollo utilizzato per migliorare e semplificare il processo di autenticazione. Alleggerisce i servizi Web che lo utilizzano dall'onere di gestire la moltitudine di credenziali degli utenti, con i meccanismi di protezione necessari che ne conseguono. I protocolli di SSO possono autenticare gli utenti mediante protocolli passwordless. Grazie a questi protocolli un utente effettua l'autenticazione tramite un token hardware in grado di eseguire operazioni di autenticazione basate su crittografia asimmetrica. Questo permette di snellire e irrobustire il processo complessivo di autenticazione. Il limite principale risiede nella centralizzazione del protocollo SSO che non è in grado di tollerare intrusioni nell'infrastruttura che lo esegue.
Infatti, come mostrano attacchi recenti \cite{volexity:solarwinds} \cite{nsa:authentication}, gli attaccanti che riescono a compromettere l'infrastruttura di SSO sono in grado di impersonare utenti arbitrari.
Per mitigare questa problematica sono nati gli studi sulla cosiddetta SSO \emph{survivability}. Tale filone si occupa di irrobustire il protocollo SSO utilizzando un approccio distribuito e replicato. Si prevede un numero di attori malevoli di cui è possibile tollerare l'intromissione, numero da cui dipende il livello di sicurezza.

L'esigenza di questa tesi nasce dal fatto che i protocolli esistenti in ambito SSO \emph{survivable} non permettono di modificare il livello di sicurezza una volta che è stato stabilito durante la fase di creazione dell'infrastruttura. Questo impedisce ai vari servizi Web di richiedere alla stessa infrastruttura SSO livelli di sicurezza adeguati alle caratteristiche del servizio offerto. Si è affrontato dunque il problema dell'integrazione del livello di sicurezza all'interno dei protocolli Passwordless e SSO preesistenti. A tal fine, si è reso necessario implementare modifiche al codice di un autenticatore open source, cui hanno fatto seguito verifiche delle performance di autenticazione e un breve studio di fattibilità per l'implementazione su hardware reale.

La tesi è strutturata nel seguente modo: nel capitolo \ref{conoscenze} vengono introdotte le conoscenze necessarie per proseguire nella lettura dei successivi; il capitolo \ref{modellazione} di modellazione in cui vengono descritti gli attori che partecipano al protocollo di SSO proposto, il funzionamento dei protocolli citati e la soluzione studiata per integrare il livello di sicurezza; il capitolo \ref{dettagli} in cui vengono approfondite le modifiche apportate al codice degli attori presenti per realizzare l'integrazione; il capitolo \ref{prestazioni} di analisi dei dati raccolti, valutazione delle tempistiche e studio finale di fattibilità; il capitolo \ref{conclusioni} conclusivo con le considerazioni finali ed eventuali sviluppi futuri.

\chapter{Modellazione del sistema}
\label{modellazione}

In questo capitolo verranno inizialmente descritti gli attori che concorrono alle fasi di registrazione/autenticazione per poi definire il flusso delle operazioni stesse.

\section{Attori}
\label{attori}

Lo schema di seguito rappresenta lo stato attuale dello standard FIDO in accordo all'implementazione descritta successivamente.
\begin{figure}[htb]
	\centering
	\includegraphics[width=.4\columnwidth]{figures/example.pdf}
	\caption{PLACEHOLDER.}
	\label{fig:esempio}
\end{figure}
% INSERIRE SCHEMA ATTORI STATO ATTUALE

\subsection{Service Provider}
\label{service_provider}

Il Service Provider è un fornitore di un generico servizio a cui l'utente è interessato ad accedere. Può essere un qualunque servizio di streaming, banking, shopping etc. il quale fa uso di un intermediario per l'autenticazione dei propri utenti. Questo può avvenire per varie ragioni sia economiche che legate alla sicurezza. 

Il Service Provider fa uso del \textbf{security level} per indicare il livello di \emph{survivability} desiderato: esso può crescere all'aumentare della confidenzialità del servizio erogato.

\subsection{Identity Provider}
\label{identity_provider}

L'Identity Provider è un ente terzo che si occupa di fornire a un Service Provider il servizio di autenticazione. Fa parte del cosiddetto \emph{Relying Party}, cioè quella parte che fornisce un accesso sicuro ad un servizio. 

Nella variante \emph{survivable}, invece che fornire un Identity Server unico con il quale il FIDO Client comunica, ne fornisce un numero \textbf{n} desiderato dal Service Provider, a seconda della sicurezza richiesta dal tal servizio. In questo modo durante la fase di creazione l'operazione dovrà essere replicata dal client su tutti gli \emph{n} Identity Server. La parte di autenticazione invece verrà svolta solo su un sottoinsieme $\mathbf{k\leq n}$ di questi. Compito dell'Identity Provider è anche quello di fornire il token di autenticazione al client da presentare al Service Provider per fare in modo che l'utente possa accedere al servizio scelto.

I vari Identity Server comunicano con il FIDO Client, tramite User Agent, per creare le credenziali degli utenti mantenendo in memoria le challenge generate e gli identificatori degli utenti registrati con le relative chiavi pubbliche ricevute. 

Compito dell'Identity Server è anche quello di rilevare eventuali tentativi di duplicazione dell'autenticatore fisico. Per fare ciò lo standard FIDO prevede un contatore gestito dall'autenticatore, sia esso globale o multiplo, che viene aggiornato ad ogni operazione avvenuta con successo. Il contatore prende il nome di \emph{signature counter}. Il server mantiene in memoria l'ultimo valore ricevuto e, all'interazione successiva, controlla che non vi siano discrepanze. In particolare se il valore ricevuto è minore di quello salvato in memoria dal server allora si può essere in presenza di un tentativo di clonazione. (L'accadere dell'inverso non è detto che costituisca un problema, motivo per il quale può essere utilizzato anche un contatore globale.) Ciò significa che un attaccante maligno ha duplicato l'autenticatore ad uno stato precedente mentre l'utilizzatore ha continuato ad autenticarsi incrementando il contatore. 

Nella variante survivable è stato modellato il sistema adottando un contatore specifico per ogni livello di sicurezza. Ad ogni autenticazione, e quindi ad ogni livello di sicurezza desiderato specificato dal Service Provider, verrà incrementato il contatore corrispondente a \emph{k}. Questa modifica non altera il funzionamento del server, che riceve sempre un valore di contatore unico, indi per cui non è stato necessario operare modifiche in questo senso dal lato del server.

\subsection{Autenticatore}
\label{autenticatore}

L'autenticatore hardware è un dispositivo che, tramite l'interazione fisica con l'utente, permette l'accesso al servizio web richiesto. Durante la fase di creazione delle credenziali, l'autenticatore si occupa di creare una coppia di chiavi crittografiche, una pubblica e una privata. Utilizzerà la chiave privata per firmare le challenge che gli vengono sottoposte mentre la pubblica verrà inviata al server così che esso possa verificare l'autenticità delle firme.

Per la fase di creazione e le successive di autenticazione viene richiesto all'utente di compiere un'azione: essa può essere la pressione di un pulsante sulla chiavetta stessa, un collegamento NFC oppure ancora l'identificazione tramite impronta digitale. Così facendo l'operazione in corso viene autenticata.

Nella variante \emph{survivable} viene inviato il digest ottenuta dall'hashing delle \emph{n} challenge dei server che partecipano all'operazione. Esso appare all'autenticatore come una challenge unica, motivo per il quale non è stato necessario modificare il codice dell'autenticatore per richiedere l'interazione ad ogni singola challenge. 

Durante la fase di registrazione viene inizializzato un contatore tramite cui vengono tenute traccia delle operazioni conseguite correttamente. Ad ogni operazione il contatore, al corrispondente livello di sicurezza, viene incrementato mono-atomicamente. Come detto in precedenza, si è reso necessario implementare un contatore specifico per ogni livello di sicurezza richiesto, sia in fase di autenticazione che di creazione. Prima di fare ciò, è stato però necessario introdurre un contatore per ogni credenziale creata, poiché allo stato attuale il codice della Solokey, per questioni di spazio, prevede un contatore globale unico. 

\subsection{Client}
\label{client}

Un dispositivo client che sfrutta un User Agent conforme ad implmentare le specifiche FIDO per il dialogo con il Relying Party e con l'autenticatore, in collaborazione con il dispositivo hardware sottostante su è installato l'User Agent, tipicamente un sistema operativo. Il client è quindi interposto tra il FIDO Server e l'autenticatore fisico, agendo da intermediario. Il suo compito è duplice:
\begin{itemize}
	\item Comunicare con il server al fine di iniziare, e successivamente terminare, le operazioni di autenticazione e di creazione delle credenziali 
	\item Comunicare con l'autenticatore allo scopo di creare le chiavi crittografiche e firmare le challenge ricevute dal server
\end{itemize}
La comunicazione è bidirezionale, duplice e segue i costrutti specificati dalle API definite nei relativi standard: CTAP per l'interazione con l'autenticatore e WebAuthn per la comunicazione con il FIDO Server. 

Nel caso particolare dell'estensione survivable il client si occupa di replicare le operazioni su \emph{n} FIDO Server distinti, computando l'hash delle challenge ricevute e fornendolo all'autenticatore come un digest unico su cui apportare la firma. 

\section{Flusso operativo}
\label{flusso_operativo}

Il flusso operativo si compone di due operazioni distinte: una fase di creazione delle credenziali e una fase di autenticazione dell'utente. Tali operazioni vengono svolte, rispettivamente, durante la registrazione al servizio del Service Provider e a tutti i login successivi.

\subsection{Fase di registrazione}
\label{registrazione}

La fase di registrazione si origina a partire dalla richiesta dell'utente, utilizzando un User Agent, di registrarsi ad un servizio, offerto da un Service Provider, che supporti l'autenticazione passworldess, tramite degli Identity Provider. Il Service Provider fornisce all'utente il livello di sicurezza \emph{n} necessario per completare l'operazione. Il processo messo in atto è il seguente:

\begin{enumerate}
	\item Ogni Identity Server crea il proprio stato interno e la challenge
	\item Ogni Identity Server invia al Client una serie di requisiti secondo cui deve essere svolta la cerimonia di registrazione, e la challenge generata
	\item Il Client salva tutte le challenge ricevute in un vettore e computa l'hash dello stesso
	\item Il Client effettua una chiamata al metodo opportuno dell'autenticatore, fornendo i requisiti di creazione richiesti dal Relying Party e il digest computato come challenge
	\item L'autenticatore procede a generare la coppia di chiavi crittografiche seguendo le imposizioni del server e invia al client la challenge firmata accompagnata dalla chiave pubblica e dal signature counter
	\item Il Client invia ad ogni Identity Server il vettore con le challenge, l'hash dello stesso, la firma. la chiave pubblica e il signature counter ricevuto
	\item Ogni Identity Server controlla che la challenge da lui generata sia presente all'interno del vettore e controlla, computando lui stesso l'hash del vettore, l'integrità di quanto ricevuto. Infine, verifica tramite la chiave pubblica fornitagli l'autenticità della firma.
	\item Qualora il processo sia andato a buon fine, gli Identity Server salveranno le informazioni ricevute (contatore, chiave pubblica, identificatore del client) al proprio interno e forniranno al Client l'attestazione con cui poter completare la registrazione presso il Service Provider
\end{enumerate} 

\subsection{Fase di autenticazione}
\label{autenticazione}

La fase di autenticazione ricalca i passaggi di quella di registrazione con la differenza che gli Identity Server sono già in possesso della chiave pubblica con cui verificare l'autenticità della firma apportata alle challenge. 

In questa fase viene comunicato un security level pari a \emph{k} da parte del Service Provider, cioè il numero di Identity Server presso cui è necessario autenticarsi. Questo valore prende in considerazione la tollerabilità alle intrusione che ha il Service Provider. In particolare: \emph{k} può essere $ \emph{(2j+1)} $, dove \emph{j} rappresenta il numero di Identity Server di cui si può tollerare la compromissione, oppure può essere $ \emph{(2j+2)} $ in caso di requisiti particolarmente stringenti. 

\begin{enumerate}
	\item Ogni Identity Server crea il proprio stato interno e la challenge
	\item Ogni Identity Server invia al Client una serie di requisiti, secondo cui deve essere svolta la cerimonia di autenticazione, e la challenge generata
	\item Il Client salva tutte le challenge ricevute in un vettore e computa l'hash dello stesso
	\item Il Client effettua una chiamata al metodo opportuno dell'autenticatore, fornendo i requisiti di autenticazione richiesti dal Relying Party e il digest computato come challenge
	\item L'autenticatore invia al Client la challenge firmata e il contatore, specifico del livello di sicurezza stabilito, aggiornato
	\item Il Client invia ad ogni Identity Server il vettore con le challenge, l'hash dello stesso, la firma ricevuta e il signature counter
	\item Ogni Identity Server controlla che la challenge da lui generata sia presente all'interno del vettore e controlla, computando lui stesso l'hash del vettore, l'integrità di quanto ricevuto. Infine, verifica tramite la chiave pubblica, memorizzata precedentemente, l'autenticità della firma
	\item Ogni Identiy Server controlla che il \emph{signature counter} ricevuto sia maggiore di quello memorizzato in precedenza e in tal caso aggiorna quest'ultimo
	\item Se i passaggi precedenti sono avvenuti con successo rilasciano al Client l'attestazione tramite cui può completare l'autenticazione presso il Service Provider
\end{enumerate}

% PAGINA VUOTA
%\clearpage\null\thispagestyle{empty}\clearpage
%\appendix
%\appendixpage
%\addappheadtotoc

%\clearpage\null\thispagestyle{empty}\clearpage


%\listoffigures


\begin{flushleft}
\bibliographystyle{plain}
\bibliography{sections/references} 
\end{flushleft}

\end{document}
